\begin{description}
\item[(a)] We consider the $\triangle ESP$, the vertices of which are Earth
  (E), Sun (S) and the minor planet (P).

% FIGURE NEEDED: triangle sketch (2020 GeCAA data analysis solutions, page 5):
% triangle E-S-P with E at top left, S at bottom, P at right; side ES labelled
% a_earth = 1 au, side EP labelled r, side SP labelled r_h; angle Delta-lambda
% at E, angle theta_h at S between ES and SP, and angle psi_P at S between SP
% and a dashed reference direction.

  where $\Delta\lambda = |\lambda - \lambda_\odot|$. \hfill(2.0 points)
  \begin{align*}
    r\ \text{(in m)} &= \frac{R_\oplus}{p\ \text{(in rad)}} \\
    \therefore\ r\ \text{(in au)} &= \frac{206265 R_\oplus}{p\cdot a_\oplus}
  \end{align*}
  \hfill(1.0 points)

  By using the cosine rule and sine rule,
  \begin{align*}
    r_h^2 &= a_\oplus^2 + r^2 - 2a_\oplus r\cos(\Delta\lambda) \\
    r_h &= \sqrt{a_\oplus^2 + r^2 - 2a_\oplus r\cos(\Delta\lambda)}
  \end{align*}
  \hfill(1.0 points)
  \begin{align*}
    \frac{\sin\theta_h}{r} &= \frac{\sin(\Delta\lambda)}{r_h} \\
    \therefore\ \theta_h &=
      \sin^{-1}\left(\frac{r\sin(\Delta\lambda)}{r_h}\right)
  \end{align*}
  \hfill(1.0 points)

  An angle $\psi_P$ is defined as angle swept by the vector
  $\overrightarrow{SP}$ from an arbitrary x axis (see figure). Let us fix the
  direction of the x-axis as the direction $\overrightarrow{SP}$ at the
  initial moment $t = 2012.3$. In other words, $\psi_P(2012.3) = 0$. Using the
  same x-axis, let us define corresponding angle for the Earth,
  $\psi_E$. Thus, we have
  \[
    \psi_P = \psi_E - \theta_h
  \]
  If we take any two consecutive measurements,
  \[
    \Delta\psi_E = \Delta\lambda_\odot
  \]
  \hfill(2.0 points)

  This is preferred over the conversion of time difference to angle as this
  data has better accuracy.

  Using these formulae, one can find $r$, $r_h$, $\psi_E$ and $\psi_P$ for
  every measurement instance in the data. $\psi_P$ and $r_h$ can be used to
  plot locations of the minor planet.

  Alternatively, one can also convert these to rectangular coordinates, i.e.
  \begin{align*}
    x &= r_h\cos\psi_P, \\
    y &= r_h\sin\psi_P;
  \end{align*}

  \begin{center}
  \begin{tabular}{rrrrrrrrr}
    \multicolumn{1}{c}{$t$} & \multicolumn{1}{c}{$r$} &
    \multicolumn{1}{c}{$r_h$} & \multicolumn{1}{c}{$\sin\theta_h$} &
    \multicolumn{1}{c}{$\theta_h$} & \multicolumn{1}{c}{$\psi_E$} &
    \multicolumn{1}{c}{$\psi_P$} & \multicolumn{1}{c}{$x$} &
    \multicolumn{1}{c}{$y$} \\
    \multicolumn{1}{c}{[year]} & \multicolumn{1}{c}{[au]} &
    \multicolumn{1}{c}{[au]} & &
    \multicolumn{1}{c}{[$^\circ$]} & \multicolumn{1}{c}{[$^\circ$]} &
    \multicolumn{1}{c}{[$^\circ$]} & \multicolumn{1}{c}{[au]} &
    \multicolumn{1}{c}{[au]} \\
    \hline
    2012.3 & 2.302 & 2.073 & $-1.0000$ & 270.00 & 270.00 & 0 & 2.073 & 0 \\
    2012.6 & 1.215 & 2.020 & $-0.4510$ & $-26.81$ & 3.88 & 30.69 & 1.737 & 1.031 \\
    2012.9 & 1.240 & 2.229 & 0.1097 & 6.30 & 111.13 & 104.83 & $-0.571$ & 2.155 \\
    2013.4 & 3.664 & 2.839 & 0.6391 & 140.27 & 293.89 & 153.62 & $-2.543$ & 1.262 \\
    2013.6 & 4.071 & 3.106 & $-0.2991$ & 197.40 & 3.64 & 166.24 & $-3.017$ & 0.739 \\
    2013.9 & 3.198 & 3.309 & $-0.9655$ & $-74.92$ & 110.87 & 185.79 & $-3.292$ & $-0.334$ \\
    2014.2 & 2.783 & 2.007 & 0.7357 & 132.63 & 222.34 & 89.71 & 0.010 & 2.007 \\
    2014.5 & 4.419 & 3.845 & 0.8735 & 119.13 & 328.09 & 208.96 & $-3.364$ & $-1.862$ \\
    2014.8 & 4.805 & 3.962 & $-0.5914$ & 216.26 & 74.50 & 218.24 & $-3.112$ & $-2.453$ \\
    2015.0 & 4.332 & 3.994 & $-0.9738$ & $-76.85$ & 149.24 & 226.09 & $-2.770$ & $-2.877$ \\
    2015.3 & 3.035 & 3.985 & 0.2736 & 15.88 & 257.6 & 241.72 & $-1.888$ & $-3.509$ \\
    \hline
  \end{tabular}
  \end{center}
  % In the original the 2014.2 row is printed in red (the outlier).

  there are 11 rows. Two points for each row. \hfill(22.0 points)

  \textbf{Note.} The parallax value for $p(2014.2) = 3.16''$ is obviously an
  outlier. Due to this, the corresponding position of the minor planet is also
  an outlier. Therefore, it should not be included in the plot or in the next
  analysis. \hfill(1.0 points)

% FIGURE NEEDED: rectangular plot (2020 GeCAA data analysis solutions, page
% 6): planet positions (blue points), visually fit ellipse (light-blue), fit
% axis (green line), approximate axis through supposed aphelion and (0,0)
% (magenta line), Sun at origin (yellow) and Earth positions (orange points),
% x and y in au. Worth 4.0 points.

  These 4 points are for merely drawing the approximate polar or rectangular
  plot. Points for the line of apsides are separate (see below). Marking the
  earth positions or drawing the visually fit ellipse (blue line) is NOT
  expected.

  One may try to identify a minimum and a maximum and maximum in the % sic: "a maximum and maximum"
  calculated values of $r_h$ and take these as the aphelion and perihelion. In
  the figure these two points are shown and the line between drawn, i.e.\ the
  magenta line. Clearly these are not the aphelion and perihelion (as
  evidenced by the fact that the line connecting them does not pass through
  the focus), but they lie close. \hfill(1.0 + 1.0 points)

  We see that the points on either side of supposed aphelion are reasonably
  close (and the planet will be moving slower close to aphelion). However,
  there is a big gap between the second and third point, we take the aphelion
  position to be the correct one and draw the line of apsides to pass through
  aphelion and (0,0). (see below for the actual calculation of $a_p$ that does
  not rely on the graphical identification of the line of the apsides.)
  \hfill(2.0 points)

\item[(b)]
  \begin{description}
  \item[(i)] Assuming that the points above are close to the perihelion and
    aphelion, we should realise that the distance between them is roughly
    equal to the length of the major axis. Using the cosine rule,
    \begin{align*}
      \psi_1 &= 30.69^\circ \\
      r_1 &= 2.020\,\mathrm{au} \\
      \psi_2 &= 226.09^\circ \\
      r_2 &= 3.994\,\mathrm{au} \\
      \therefore\ \Delta\psi_P &= 226.09^\circ - 30.69^\circ \\
        &= 195.40^\circ \\
      2a_P &\approx \sqrt{r_1^2 + r_2^2 - 2r_1 r_2\cos\Delta\psi_P} \\
        &= \sqrt{2.020^2 + 3.994^2
           - 2\times2.02\times3.994\times\cos195.40^\circ} \\
        &= 5.965\,\mathrm{au} \\
      \therefore\ a_P &\approx 2.98\,\mathrm{au}
    \end{align*}
    \hfill(3.0 points)

  \item[(ii)]
    \[
      e = \frac{r_a - r_p}{r_a + r_p}
        \approx \frac{3.994 - 2.02}{3.994 + 2.02}
        \approx \frac{1}{3} = 0.33
    \]
    \hfill(2.0 points)

    Some students may visually fit an ellipse to the data and draw this line
    up to this visually fit ellipse. For the fit ellipse in the figure, the
    length of the major axis (green line) is $2a = 2.98$\,au and
    $e = \frac{1}{3}$, rotated by $52^\circ$. % sic: 2a = 2.98 au (cf. 2a_P = 5.965 au above)

  \item[(iii)] From Kepler's Laws,
    \begin{align*}
      P &= \sqrt{a_P^3} = \sqrt{2.98^3} \\
        &= 5.15\,\mathrm{yr}
    \end{align*}
    \hfill(1.0 points)

    The semi-period will be 2.58\,yr. Note that the time difference between
    our chosen points is only 2.4\,yr. This is fine, as we can see from the
    best fit ellipse that these points are not exactly along the line of
    apsides.

    At the same time, this is precisely the reason why one should not take
    $\Delta t$ between these two points as the semi-period.
  \end{description}

\item[(c)] The error in the semimajor axis is determined by the error in $r_1$
  and $r_2$ which, in turn, depend on $p$ and $\Delta\psi$. The latter is
  dominated by the uncertainty in the orientation of the ellipse (i.e., the
  assumption that $r_1$ and $r_2$ are close to the line of the apsides).

  In order to take this into account we take as $\Delta\psi_P$ the difference
  with the value that corresponds to the difference in angles of the actual
  perihelion and aphelion ($180^\circ$); i.e.,
  $\delta\Delta\psi_P \approx 15/180$, corresponding to
  $\Delta(\Delta\psi) \approx 0.08\pi \approx 0.26$. The geocentric distance
  is determined via parallax, the other quantities are constants, so it may be
  written $r_g \propto p^{-1}$. The immediate consequence is
  $\delta r_g = \delta p$ where $\delta$ is the designation for the relative
  error. Thus,
  \[
    \Delta r = r\delta r = r\delta p
  \]
  This gives $\Delta r_1 = 1.01\times10^{-2}$\,au and
  $\Delta r_2 = 1.997\times10^{-2}$\,au.

  The error in $a_p$ is then
  \begin{align*}
    \Delta a_P &= \frac{\sqrt{(r_1 - r_2\cos\Delta\Psi_p)^2\Delta r_1^2
      + (r_2 - r_1\cos\Delta\Psi_p)^2\Delta r_2^2
      + (r_1 r_2\sin\Delta\Psi_p)^2\Delta(\Delta\Psi_p)^2}}{2\sqrt{2}a_P} \\
      &= 0.04\,\mathrm{au} \\
    \therefore\ a_P &= (2.98\pm0.04)\,\mathrm{yr} % sic: units given as yr in source, should be au
  \end{align*}
  \hfill(3.0 points)

  The uncertainty in the period can then be obtained from Kepler's law and is
  \begin{align*}
    \delta P &= 3/2\,\delta a_p = 2\times10^{-2} \\
    \therefore\ P &= (5.2\pm0.1)\,\mathrm{yr}
  \end{align*}
  \hfill(1.0 points)

  The eccentricity can be calculated from $r_p = a_p(1 - e)$. Error
  propagation leads to
  \begin{align*}
    \Delta e &= \sqrt{\left((1 - e)\delta r\right)^2
      + \left(\frac{1 - e}{a_p}\Delta a_p\right)^2} \\
      &= 9\times10^{-3} \approx 0.01 \\
    \therefore\ e &= 0.33\pm0.01
  \end{align*}
  \hfill(2.0 points)
\end{description}
